% Part: normal-modal-logic
% Chapter: frame-definability
% Section: first-order-definability

\documentclass[../../../include/open-logic-section]{subfiles}

\begin{document}

\olfileid{nml}{frd}{fol}

\olsection{Definabilitas Orde Kapisan}

Kita wis weruh yen saweneh sipat relasi aksesibilitas bingkai
bisa didéfinisikake dening !!{formula}s modal. Upamane, simetri
bingkai bisa didéfinisikake dening !!{formula}~\Ax{B}, $p \lif
\Box\Diamond p$. Kabeh kondisi kang wis kita temoni nganti saiki
bisa diandharake nganggo !!{formula}s orde kapisan ing !!a{language}
kang mung ngemot siji !!{predicate} rong panggonan. Upamane, simetri
didéfinisikake dening
$\lforall[x][\lforall[y][(\Atom{Q}{x,y} \lif \Atom{Q}{y, x})]]$
ing pangertèn yen !!{structure} orde kapisan~$\Struct{M}$ kanthi
$\Domain{M} = W$ lan $\Assign{Q}{M} = R$ nyukupi formula kasebut yen
lan mung yen $R$ simetris. Iki nuntun marang definisi iki:

\begin{defn}
  Kelas bingkai~$\mClass{F}$ \emph{bisa didéfinisikake kanthi orde
    kapisan} yen ana !!a{sentence}~$!A$ ing basa orde kapisan kang mung
  nduweni siji !!{predicate} rong panggonan~$Q$ saengga $\mModel{F} =
  \tuple{W, R} \in \mClass{F}$ yen lan mung yen $\Sat{M}{!A}$ ing
  !!{structure} orde kapisan~$\Struct{M}$ kanthi $\Domain{M} = W$ lan
  $\Assign{Q}{M} = R$.
\end{defn}

Nyatane sipat-sipat lan !!{formula}s modal kang
ndéfinisikake sipat kasebut, kang wis ditliti nganti saiki, iku
istimewa. Ora saben !!{formula} ndéfinisikake kelas bingkai kang bisa
didéfinisikake kanthi orde kapisan; lan ora saben kelas bingkai kang
bisa didéfinisikake kanthi orde kapisan bisa didéfinisikake dening
formula modal.

Tuladha kontra kanggo pratelan kapisan yaiku formula L\"ob:
\begin{equation}
\Box(\Box p \lif p) \lif \Box p. \tag{\Ax{W}}
\end{equation}
\Ax{W} ndéfinisikake kelas bingkai transitif kang mapan kanthi becik
saka arah kosok baline. Sawijining relasi mapan kanthi becik yen ora
ana runtutan tanpa wates
$w_1$, $w_2$, \dots{} kang nyukupi $Rw_2w_1$, $Rw_3w_2$, \dots.
Upamane, relasi $<$ ing~$\Nat$ mapan kanthi becik, nanging relasi
$<$ ing~$\Int$ ora. Relasi mapan kanthi becik saka arah kosok baline
yen lan mung yen relasi kosok baline mapan kanthi becik. Dadi, relasi
kang mapan kanthi becik saka arah kosok baline yaiku relasi kang ora
nduweni runtutan tanpa wates $w_1$, $w_2$, \dots{} kang nyukupi
$Rw_1w_2$, $Rw_2w_3$, \dots.

Nanging ora ana !!{formula} orde kapisan kang
ndéfinisikake relasi transitif kang mapan kanthi becik saka arah
kosok baline. Upamane~$\Sat{M}{!F}$ yen lan mung yen $R =
\Assign{Q}{M}$ transitif lan mapan kanthi becik saka arah kosok
baline. Kanggo saben indeks wiwit loro, $!A_n$ iku !!{formula}
\[
(\Atom{Q}{a_1,a_2} \land \dots \land
    \Atom{Q}{a_{n-1},a_{n}})
\]
Saiki delengen himpunan !!{formula}s
\[
\Gamma = \{!F, !A_2, !A_3, \dots\}.
\]
Saben subhimpunan winates saka $\Gamma$ bisa dipenuhi: Yen ora ana
formula rantai ing subhimpunan iku, pilih model kanthi siji donya lan
relasi kosong. Yen ana, pilih $k$ paling gedhe saengga $!A_k$ kalebu ing
subhimpunan iku, lan pilih $\Domain{M_k} = \{1, \dots, k\}$,
$\Assign{a_i}{M_k} = i$ kanggo indeks nganti k (konstanta liyane
ditafsirake sakarepe), lan $\Assign{Q}{M_k} = <$. Amarga $<$ ing
$\{1, \dots, k\}$ transitif lan mapan kanthi becik saka arah kosok
baline, $\Sat{M_k}{!F}$. Miturut konstruksi, $\Sat{M_k}{!A_i}$ kanggo
saben $i \le k$. Miturut Teorema Kompaksi kanggo logika orde kapisan,
$\Gamma$ bisa dipenuhi ing sawijining !!{structure}~$\Struct{M}$.
Miturut hipotesis, amarga $\Sat{M}{!F}$, relasi $\Assign{Q}{M}$ mapan
kanthi becik saka arah kosok baline. Nanging cetha yen
$\Assign{a_1}{M}$, $\Assign{a_2}{M}$, \dots{} bakal mbentuk runtutan
tanpa wates saka jinis kang dilarang dening sipat kasebut.

Tuladha kontra kanggo pratelan kapindho yaiku sipat
universalitas: kanggo saben $u$ lan $v$, $Ruv$. Bingkai universal bisa
didéfinisikake kanthi orde kapisan dening formula
$\lforall[x][\lforall[y][\Atom{Q}{x,y}]]$. Nanging ora ana formula
modal kang valid ing kabeh lan mung bingkai universal. Iki akibat saka
asil kang uga narik kawigaten dhewe: formula-formula kang valid ing
bingkai universal persis padha karo kang valid ing bingkai refleksif,
simetris, lan transitif. Ana bingkai refleksif, simetris, lan
transitif kang ora universal; mula saben !!{formula} kang valid ing
kabeh bingkai universal uga valid ing sawetara bingkai kang ora
universal.

\end{document}
